\label{sec:introduction}

When the Maricopa County Citizens Redistricting Commission drew Arizona's
legislative districts in 2011, it held dozens of public hearings at which
residents testified about the economic communities that deserved to be kept
whole in a single district.
Ranchers and farmers in Yuma testified that their water rights, commodity
markets, and co-op memberships made them an economic community distinct from
suburban Tucson.
Autoworkers in Maricopa County testified that their union locals, plant
bargaining units, and supplier relationships made them a community distinct
from adjacent white-collar suburbs.
The Commission listened, noted the testimony, and then---lacking any
computational tool to operationalize economic community membership---drew
lines based primarily on population equality and VRA compliance.
The economic communities were an aspiration, not an input.

This paper provides the computational input.

The M-track community character framework, introduced in M.0
\citep{m0framework}, operationalizes diverse community character signals
as edge weights in the census-tract adjacency graph.
Communities that belong together get heavy edges; communities that are
naturally distinct get light edges.
METIS's minimum-cut objective then respects community structure automatically:
cuts run along light edges (natural community transitions) and avoid heavy
edges (strong community bonds).

M.1 is the first M-track implementation paper.
It focuses on \emph{economic character}: the mix of commercial, industrial,
and residential activity in a census tract, as measured by the Census
Bureau's LODES Workplace Area Characteristics (WAC) data
\citep{abowd2009lodes}.
Economic character is the right first M-track signal because:

\begin{enumerate}
  \item \textbf{Data quality.} LODES WAC is derived from administrative
    payroll records (Unemployment Insurance wage records) with near-complete
    employer coverage, not survey estimates with sampling error.
  \item \textbf{Legal recognition.} Economic communities of interest are
    explicitly recognized in redistricting criteria in California, Colorado,
    Arizona, Washington, and multiple other states
    \citep{caprop11,coart5s44,azircprop106,warcw44}.
  \item \textbf{Partisan neutrality.} No electoral, racial, or demographic
    data enters the computation.
    The signal is the economic profile of the workplace geography.
  \item \textbf{Empirical effect.} As this paper demonstrates, the signal
    produces a substantial proportionality improvement in NC-14, the largest
    improvement from any Layer~2 modification we have tested on that state.
\end{enumerate}

\paragraph{Distinction from B.27.}
The companion paper B.27 \citep{b27econchar} presents the same algorithm
from an implementation and benchmarking perspective: data pipeline, Rust code
structure, test invariants, computational performance, and comparative analysis
of the three test states.
M.1 presents the same algorithm from the community character perspective:
which economic communities does it preserve, how does it fit the M-track
framework, what is the primary M-track metric (within-district economic
variance), and how is it legally deployed.

\paragraph{Paper organization.}
Section~\ref{sec:background} describes the LODES WAC data source and the
three economic character signals.
Section~\ref{sec:framework} defines the character vectors, cosine similarity,
and edge weight formula within the M-track framework.
Section~\ref{sec:results} presents empirical results including the
proportionality gap comparison and a discussion of within-district economic
variance.
Section~\ref{sec:legal} provides a detailed legal usage guide for redistricting
practitioners.
Section~\ref{sec:conclusion} concludes.
